\documentclass{jsaishoshiki}
%
%  初版 2026/03/16
%  第２版 2026/03/19 空白調整
%  第３版 2026/03/30
%
\title{論文題名}
\author{〇情報　太郎\supscript{1)},農業　花子\supscript{2)}}
\affiliation{1)独立行政法人農業情報研究所，〒123-4567　茨城県つくば市中央890 \\
2)農業情報大学大学院情報研究科，〒987-6543　東京都情報区農業町210 }

\begin{document}

\maketitle

\begin{abstract}
  \vspace{1cm}
  
  要旨は，500文字程度で作成してください．要旨本体は，本文のフォント設定と同じ「明朝体（Times New Roman）」「9ポイント」です．
  \\
  \\
  \\
\end{abstract}

\keywords{キーワード１，キーワード２，キーワード３，キーワード４，キーワード５（最大５つまで）}

\begin{multicols}{2}
  
\section{緒言（大見出し, 11P）}
「緒言」の２つ上の行までが第１セクション，「緒言」の１つ上の行以降が第２セクションになります．

第２セクションには本文（図表を含む）と引用文献が含まれます．

第１セクションは１段組，第２セクションは２段組です．

\section{ページ数について}

一般講演およびオーガナイズドセッションについては2ページです．シンポジウムについては，2ページまたは4ページです．


\section{タイトルからキーワードまで（大見出し）}
\subsection{タイトル・著者氏名・所属について（小見出し）}

タイトルは「ゴシック」「15ポイント」「センタリング」として下さい．

タイトルの後に１行空白（標準：９ポイント）を設定して，著者氏名を明記して下さい．著者氏名は「明朝体」「13ポイント」「センタリング」として下さい．

著者氏名の後に１行空白（標準：９ポイント）を設定して，所属を明記して下さい．所属は「明朝体」「9ポイント」「センタリング」として下さい．

\subsection{要旨とキーワードについて（小見出し）}

要旨は500文字程度で作成して下さい．「要旨」とある部分は「ゴシック体」「11ポイント」です．要旨本体は，後述する本文のフォント設定と同じです（明朝体9ポイント）．

キーワードは最大5つ作成して下さい．「キーワード」とある部分は「ゴシック体」「11ポイント」です．キーワード本体は，後述する本文のフォント設定と同じです（明朝体9ポイント）．

\section{本文について（大見出し）}
\subsection{本文のフォントについて（小見出し）}

本文のフォントはサイズが「９ポイント」，種類は和文が「明朝体」，欧文が「Times New Roman」です．数字はアラビア数字（1,2,3,など）で１桁，２桁ともに半角（Times New Roman）です．

\subsubsection{見出しについて（小見出し,9P）}
「緒言」をはじめとした大見出しは「ゴシック」「11ポイント」です．上下に１行（標準：９ポイント）の空白行を設定します．
              
小見出しは「ゴシック」「９ポイント」です．小見出しの上に1行（標準：９ポイント）の空白行を設定します．ただし，小見出しの上に大見出しがある場合は，大見出しと小見出しの間の空白行は１行（標準：９ポイント）とします．


\section{図表のタイトルについて（大見出し）}
\subsection{図のタイトルについて（小見出し）}
図のタイトルは，「図1　図のタイトル」の形式として，図の下方に設置して下さい．図番号は半角として下さい．

\subsection{表のタイトルについて（小見出しレベル）}
表のタイトルは「表1　表のタイトル」の形式として，表の上方に設置して下さい．表番号は半角として下さい．


\begin{figure}[H]
  \centering
  \includegraphics[width=65mm,height=40mm]{example-image-a}
  \caption{図のタイトルについて（小見出し）}
  \label{fig:layout_1}
\end{figure}

\section{謝辞について（大見出し）}
謝辞が必要な場合は，本文と引用文献の間に設定して下さい．見出しの「謝辞」は大見出しと同じ書式（ゴシック体11ポイント），謝辞本体は本文と同じ書式（明朝体9ポイント）で作成して下さい．

\section{引用文献について（大見出しレベル）}
引用文献については，見出しの「引用文献」は大見出しと同じ書式の「ゴシック体」「11ポイント」です．各文献については，8ポイントで作成して下さい．

文献は和文・欧文混在として，著者氏名でアルファベット順に並べて下さい．

引用例）

  山田・石田（1991）は農業におけるパーソナルコンピュータ利用の将来を展望した．一方，ハードウエア価格の急速な低下は農村において予測をはるかに上回る普及をもたらしている（宮沢 1989，伊藤ら 1992）．その傾向は海外でも同様で（Smith et al． 1991）．．．

一方，開発技術の革新（Thomas 1999）により農業向けアプリケーションは病害虫防除ソフト（Nimiya and Tsukuba 1998，農業インターネットセンター <http://www.agic.ne.jp/padb/> 1999）や経営診断ソフト（山下 2002），水稲生育診断ソフト（<http://agrinfo.narc.go.jp>，2003年1月10日参照）など続々と開発されている（Macrosoft Inc. 2002）．．．．


\section{謝辞}
本報告は，○○補助金による研究成果に基づく．

\section{引用文献}
伊藤　稔・二宮正士・遠藤　勲 ら（1992）ハード価格の推移に関する分析，農業情報研究，1:651-670．

Macrosoft Inc.(2002) Software review report, <http://www.macrosofts.com/>, browsed on Dec. 5, 2002.

宮沢喜一（1989）育種・栽培現場における小型計算機普及事例研究，育雑，36:251-256．

Nomiya, S and T. Tsukuba (1998) IMP software revolution, Agric. Info. Res., 7:11-14.

農業インターネットセンター（1999）病害虫防除支援システム，<http://www.agic.ne.jp/padb/>，2002年10月1日参照．

Smith, A., Y．Yamada and T．Sato (1991) Utilization of personal computers in agriculture, J. Agr. Info. Sci., 23: 25-32．

Thomas, J. (1999) Java RMI implementation, ed. Laurenson, M., 'Java Development Handbook', NARC Press, Tsukuba, 321-340.

山下五郎（2002）農業経営診断ソフトのすべて，JSAI編，「農業情報年鑑2003」，電農出版，東京，125-136．

山田太郎・石田光治（1991）「農業におけるパーソナルコンピュータ利用」，農業社，土浦，255pp．

\end{multicols}

\end{document}

